\documentclass[11pt, a4paper]{article}
\usepackage[utf8]{inputenc}
\usepackage[T1]{fontenc}
\usepackage{lmodern}
\usepackage[spanish, es-tabla]{babel}
\usepackage[margin=1.8cm]{geometry}
\usepackage{amsmath, amssymb}
\usepackage{xcolor}
\usepackage[most]{tcolorbox}
\usepackage{booktabs}
\usepackage{enumitem}
\usepackage[american]{circuitikz}

\setlength{\parindent}{0pt}
\setlength{\parskip}{0.5em}

% Definición de colores
\definecolor{ucbblue}{RGB}{0, 51, 102}
\definecolor{ucborange}{RGB}{204, 102, 0}
\definecolor{resultgreen}{RGB}{230, 245, 230}
\definecolor{resultborder}{RGB}{34, 139, 34}
\definecolor{warningred}{RGB}{255, 230, 230}
\definecolor{warningborder}{RGB}{200, 0, 0}

% Entornos personalizados
\newtcolorbox{expected}[1][]{
    colback=resultgreen,
    colframe=resultborder,
    fonttitle=\bfseries,
    title=Resultado Esperado #1,
    boxrule=0.8pt,
    left=5pt, right=5pt, top=3pt, bottom=3pt
}

\newtcolorbox{safety}[1][]{
    colback=warningred,
    colframe=warningborder,
    fonttitle=\bfseries,
    title=\textbf{Nota de Seguridad} #1,
    boxrule=0.8pt,
    left=5pt, right=5pt, top=3pt, bottom=3pt,
    icon=\faExclamationTriangle % Requiere fontawesome, mejor lo omitimos si no esta
}

\newcommand{\phase}[1]{\vspace{0.4cm}\noindent\colorbox{ucbblue}{\textcolor{white}{\textbf{\large #1}}}\vspace{0.2cm}}
\newcommand{\subphase}[1]{\vspace{0.2cm}\noindent\textbf{\textcolor{ucbblue}{#1}}\vspace{0.1cm}}

\begin{document}

\begin{center}
    {\LARGE \textbf{Protocolo Oficial de Pruebas de Laboratorio}}\\[0.2cm]
    {\Large \textbf{Hito 1 — Circuitos Electrónicos III (IMT-221)}}\\[0.3cm]
    \hrule
\end{center}
\vspace{0.4cm}

Este documento establece el procedimiento estandarizado para la validación física del hardware del Hito 1. Proceda únicamente con la alternativa que corresponda a la topología asignada a su grupo.

% ======================================================================
% ALTERNATIVA A: MECATRÓNICA
% ======================================================================
\section*{\textcolor{ucborange}{Alternativa A: Topología Mecatrónica}}

\phase{Fase 1: Verificación de la Red de Recorte (Clamping) del ADC}
Esta fase se realiza con el motor completamente apagado y desconectado de la red de protección. El objetivo es certificar que la barrera de diodos limite correctamente la tensión antes de integrarla al sistema de potencia.

\subphase{1. Configuración del Hardware:}
\begin{itemize}[leftmargin=*]
    \item Alimenta la red de redundancia conectando el cátodo del diodo 1N4148 ($D_{\text{FAST\_2}}$) al riel de +5 V de la fuente de alimentación (o al pin de 5V del microcontrolador).
    \item Asegúrate de que la resistencia de carga $R_{\text{LOAD}}$ ($10\,\text{k}\Omega$) esté conectada en paralelo con el diodo Zener en el nodo de salida hacia el ADC.
    \item Interconecta la tierra del generador de funciones, la tierra de la fuente de poder y las tierras de las sondas del osciloscopio en el GND común del protoboard.
\end{itemize}

\subphase{2. Inyección de la Señal de Falla:}
\begin{itemize}[leftmargin=*]
    \item Enciende el Generador de Funciones (con la salida apagada o \textit{output off}).
    \item Configura una señal senoidal de $15\,\text{Vpp}$, a una frecuencia de $1\,\text{kHz}$, con un offset de CC de $+2.5\,\text{V}$. (Esta señal oscilará críticamente entre $-5\,\text{V}$ y $+10\,\text{V}$, simulando una sobretensión masiva de falla).
    \item Conecta el terminal positivo del generador directamente en el extremo de entrada de la resistencia serie $R_S$ (el valor calculado por los estudiantes, ej. $180\,\Omega$ o $220\,\Omega$).
\end{itemize}

\subphase{3. Medición y Captura:}
\begin{itemize}[leftmargin=*]
    \item Conecta la sonda del Canal 1 del osciloscopio en el punto de inyección del generador (entrada cruda).
    \item Conecta la sonda del Canal 2 en el nodo de salida protegido (en bornes de la resistencia de carga $R_L$ o entrada al ADC).
    \item Ajusta ambos canales del osciloscopio en acoplamiento CC (\textit{DC Coupling}) y activa la salida del generador.
\end{itemize}

\begin{expected}
En la pantalla verás la senoidal completa de $15\,\text{Vpp}$ en el Canal 1, mientras que el Canal 2 mostrará una señal perfectamente recortada con un techo plano en $\approx 5.1\,\text{V}$ (acción del Zener BZX55C5V1 y el 1N4148 de redundancia) y un piso plano en $\approx -0.7\,\text{V}$ (acción del 1N4148 de piso). Exporta los datos en un archivo \texttt{.csv} para que el equipo calcule el RMSE en su Jupyter Notebook.
\end{expected}

\vspace{0.2cm}
\phase{Fase 2: Verificación de la Planta de Potencia y Supresión Inductiva}
Esta fase valida el comportamiento dinámico del motor y la efectividad del diodo Flyback. El generador de funciones se desconecta por completo.

\subphase{1. Configuración del Hardware:}
\begin{itemize}[leftmargin=*]
    \item Conecta un extremo del motorreductor amarillo a la salida de $+12\,\text{V}$ de la fuente de alimentación.
    \item Conecta el otro extremo del motor en serie a la resistencia $R_{\text{SHUNT}}$ ($1\,\text{k}\Omega$), y el extremo restante del Shunt a GND.
\end{itemize}
\begin{tcolorbox}[colback=warningred, colframe=warningborder, title=\textbf{Nota de Seguridad:}]
Para proteger los componentes ante un posible error de conexión, limita estricta y manualmente la corriente de la fuente de alimentación a un máximo de \textbf{100 mA}.
\end{tcolorbox}

\subphase{2. Captura del Transitorio SIN Flyback (Pico Inductivo Destructivo):}
\begin{itemize}[leftmargin=*]
    \item Retira temporalmente el diodo 1N4007 del motor.
    \item Conecta la sonda del Canal 1 directamente en el terminal del motor que se conecta a los 12V.
    \item Configura el trigger del osciloscopio en modo Single-Shot, flanco de bajada (\textit{falling edge}) y ajusta el nivel de disparo a $6\,\text{V}$.
    \item Realice una conmutación manual rápida: desconecte físicamente el cable de 12V del motor del riel de alimentación y vuélvalo a conectar rápidamente (\textit{tapping}).
\end{itemize}

\begin{expected}
El osciloscopio congelará en pantalla un pico transitorio negativo/positivo masivo de alta tensión (que puede superar los $100\,\text{V}$ en una fracción de microsegundo) provocado por el colapso del campo magnético de la bobina del motor.
\end{expected}

\subphase{3. Captura del Transitorio CON Flyback (Sujeción de Seguridad):}
\begin{itemize}[leftmargin=*]
    \item Vuelva a conectar el diodo 1N4007 en paralelo directo a los terminales del motor. ¡Atención a la polaridad! El cátodo (banda plateada) va al terminal de $+12\,\text{V}$ y el ánodo al terminal hacia la Shunt/GND.
    \item Repita la conmutación manual rápida (\textit{tapping}).
\end{itemize}

\begin{expected}
El osciloscopio capturará que el pico inductivo destructivo ha desaparecido por completo. La tensión en bornes de la bobina queda ahora sujeta de forma segura a un máximo de $12.7\,\text{V}$ ($V_{CC} + V_{\text{on\_diodo}}$) gracias a la malla de rueda libre.
\end{expected}

\vspace{0.2cm}
\phase{Fase 3: Integración Total y Simulación de Fallas de Corriente}

\subphase{1. Configuración del Hardware:}
Mantén el lazo del motor conectado a la fuente de $12\,\text{V}$ (con el diodo Flyback 1N4007 instalado) en serie con la resistencia $R_{\text{SHUNT}}$ ($1\,\text{k}\Omega$) a tierra. Conecta un cable desde el nodo superior de la Shunt hacia la entrada de la resistencia serie de protección $R_S$ probada en la Fase 1.

\subphase{2. Prueba en Funcionamiento Normal:}
Energiza la planta a $12\,\text{V}$ (el motor girará continuamente). Mide la tensión en el Shunt (Canal 1) y en la entrada al ADC (Canal 2).
\begin{expected}
La corriente en vacío es muy baja; la caída en el Shunt será pequeña ($\approx\,\text{mV}$ o $<2\,\text{V}$). El Canal 1 y 2 mostrarán la misma señal de CC limpia, ya que los diodos protectores están apagados (OFF).
\end{expected}

\subphase{3. Simulación de Falla por Sobretensión:}
Retira momentáneamente el cable de la Shunt y simula una sobretensión (rotor trabado o cortocircuito) inyectando una rampa de CC ajustable de $0$ a $10\,\text{V}$ desde la fuente directamente en el nodo de entrada de la Shunt.
\begin{expected}
La tensión de entrada (Canal 1) sube libremente hasta $10\,\text{V}$. En el Canal 2 (salida al ADC), el voltaje sube linealmente y se engancha exactamente en $\approx 5.1\,\text{V}$ en cuanto la entrada supera ese nivel.
\end{expected}

\newpage
% ======================================================================
% ALTERNATIVA B: BIOMÉDICA
% ======================================================================
\section*{\textcolor{ucborange}{Alternativa B: Topología Biomédica}}

\phase{Fase 1: Verificación de la Red de Recorte (Falla de Alta Tensión)}
Esta prueba valida la capacidad de la barrera de diodos para sujetar sobretensiones masivas, protegiendo el pin analógico (ADC) antes de conectar el microcontrolador.

\subphase{1. Configuración y Señal de Falla:}
\begin{itemize}[leftmargin=*]
    \item Conecta el cátodo del diodo rápido 1N4148 al riel de $+5\,\text{V}$. Coloca $R_{\text{LOAD}} = 220\,\Omega$ a la salida.
    \item Configura el Generador a $15\,\text{Vpp}$, $1\,\text{kHz}$, offset $+2.5\,\text{V}$. Inyecta en el divisor resistivo ($R_{\text{DIV}} = 10\,\text{k}\Omega$).
\end{itemize}

\begin{expected}
(Canal 1 en la entrada cruda, Canal 2 en el ADC). Canal 1 oscila entre $-5\,\text{V}$ y $+10\,\text{V}$. Canal 2 recortado perfectamente en $\approx 5.1\,\text{V}$ (Zener) y $\approx -0.7\,\text{V}$ (1N4148 de piso). Exporte a CSV.
\end{expected}

\phase{Fase 2: Simulación de Ruido por Inducción (EMI) y Validación Flyback}
Demuestra cómo un actuador clínico ruidoso acopla ruido en la línea del biosensor, y cómo el Flyback lo suprime.

\subphase{1. Configuración del Hardware y Captura de Interferencia:}
\begin{itemize}[leftmargin=*]
    \item Aísle eléctricamente la etapa del motor ($12\,\text{V}$) del biosensor. Acerque el cable de entrada del biosensor a los cables del motor para propiciar acoplamiento electromagnético.
    \item \textbf{Sin Flyback:} Retira el 1N4007. Pon osciloscopio en Peak Detect (AC Coupling). Haz \textit{tapping} con el cable del motor a 12V. Aparecerán ráfagas de espigas de ruido en la línea del biosensor.
    \item \textbf{Con Flyback:} Conecta el 1N4007 en paralelo al motor. Repite el \textit{tapping}.
\end{itemize}

\begin{expected}
Las espigas de ruido en la sensible línea analógica del biosensor disminuirán drásticamente de amplitud, demostrando la supresión de EMI en la fuente.
\end{expected}

\phase{Fase 3: Medición de Corriente de Fuga (Norma IEC 60601)}
Certifica que, ante una falla interna del MCU, $R_S$ limita la corriente hacia el paciente a un nivel seguro.

\begin{center}
\begin{circuitikz}[american, scale=1.1, transform shape]
    % Fallas y MCU
    \draw (6,3) node[vcc]{\textbf{+5V (Falla Interna)}} -- (6,1.5);
    \draw (6,1.5) to[short, -*] (6,0) node[above]{\textbf{Nodo ADC}};
    
    % Diodos de protección
    \draw (6,0) -- (6,-1) to[zzD, l_=$D_Z$] (6,-3) node[ground]{};
    \draw (7.5,0) -- (7.5,-1) to[D, l=$1N4148$] (7.5,-3) node[ground]{};
    \draw (6,0) -- (7.5,0);
    \node at (6.75, -2.5) [right] {OFF};

    % Resistencia serie y medición
    \draw (6,0) to[R, l_=$R_S$ (Protección), i_=$I_{\text{fuga}}$] (0,0);
    \draw (0,0) node[circ]{} node[left]{\textbf{Electrodo (Paciente)}};
    \draw (0,0) to[ammeter, l=Multímetro ($\mu$A)] (0,-3) node[ground]{};
\end{circuitikz}
\end{center}

\subphase{1. Configuración del Hardware y Medición:}
\begin{itemize}[leftmargin=*]
    \item Simula la falla catastrófica del MCU conectando un cable desde $+5\,\text{V}$ directo al nodo del ADC (extremo derecho de $R_S$).
    \item Conecta un multímetro digital (escala $\mu\text{A}$ CC) en serie entre la entrada del biosensor (extremo izquierdo de $R_S$) y GND.
\end{itemize}

\begin{expected}
Los diodos están en inversa para este nivel DC. Toda la corriente pasa por $R_S$. Los estudiantes verificarán que $I_{\text{fuga}} = \frac{5\,\text{V}}{R_S} \le I_{\text{límite\_seguro}}$, cumpliendo la norma IEC 60601 para partes aplicadas B/BF. Apaga la fuente y retira el puente al terminar.
\end{expected}

\end{document}
